A pesar de que el ICFES administra de forma sistemática las pruebas Saber 11 y Saber Pro, y de que ambos instrumentos evalúan competencias comparables en dos momentos del proceso formativo, la información que generan permanece desarticulada: los resultados de cada prueba se almacenan y analizan de manera independiente, sin un mecanismo que vincule el desempeño inicial de un estudiante con su desempeño final \cite{bogoya2013modelo}. La falta de integración entre estas fuentes de información dificulta comprender hasta qué punto los resultados obtenidos en la prueba Saber Pro pueden explicarse a partir de las características académicas y socioeconómicas con las que los estudiantes ingresan a la educación superior \cite{bogoya2017value}. Como consecuencia, resulta complejo diferenciar el aporte que realiza la institución durante el proceso formativo, es decir, el valor agregado \cite{mendoza2026valor}, de aquellos factores que ya estaban presentes antes de iniciar la formación universitaria. 

En la Universidad de San Buenaventura, esta situación limita la disponibilidad de herramientas que permitan responder, con evidencia cuantitativa, preguntas fundamentales para la gestión académica. Por ejemplo, cuál podría ser el desempeño esperado de un estudiante en Saber Pro de acuerdo con sus resultados en Saber 11, qué factores del contexto influyen en esa evolución \cite{rodriguez2022valor} y hasta qué punto los resultados obtenidos terminan siendo mejores o peores de lo que inicialmente se proyectaba. Sin un sistema que modele explícitamente esta relación de entrada y salida, las decisiones sobre acompañamiento estudiantil, fortalecimiento curricular y asignación de recursos continúan tomándose sin una base predictiva sólida.

De lo anterior se desprende el problema que orienta esta investigación: la inexistencia, en la Universidad de San Buenaventura, de un sistema de información capaz de integrar los resultados de las pruebas Saber 11 y Saber Pro y de modelar predictivamente la evolución de las competencias académicas entre ambos momentos. En consecuencia, el presente trabajo se formula en torno a la siguiente pregunta de investigación:

¿En qué medida es posible estimar el desempeño de los estudiantes de la Universidad de San Buenaventura en las pruebas Saber Pro a partir de sus resultados en Saber 11 y de variables socioeconómicas y de contexto, de modo que el sistema permita cuantificar la evolución de sus competencias académicas?